% Source: notes/2026-10-22-relation-to-learning.jpg (note 11)
%         notes/2026-10-22-mc-as-agnostic-learning.jpg (notes 12-13)
%         notes/2026-10-22-weak-agnostic-learning.jpg (notes 14-15)
%         notes/2026-10-22-auditing.jpg (notes 16-17)

\section{Multigroup fairness and learning}

\begin{frame}{Relation to learning}
  We've viewed groups $\G$ as subgroups of the population
  \textcolor{muted}{(potentially overlapping).}

  Equally, think of them as a class of functions
  $\C$, \ $c : \X \to \{0,1\}$.

  \pause
  \vspace{0.8em}
  Suppose $p$ is \alert{multiaccurate} for $\C$:
  \[
    \E\big[\,c(x)\,(y - p(x))\,\big] \approx 0 \qquad \forall c \in \C .
  \]

  An \alert{indistinguishability} condition: no function in $\C$ is
  correlated with the residual $y - p(x)$.
\end{frame}

\begin{frame}{For multicalibration}
  Suppose $p : \X \to [0,1]$ is \alert{multicalibrated} for $\C$:
  \[
    \E\big[\,c(x)\,(y - p(x)) \mid p(x) = v\,\big] \approx 0
    \qquad \forall c \in \C,\ \forall v .
  \]
\end{frame}

\begin{frame}{Multicalibration as agnostic learning}
  \begin{columns}[c]
    \begin{column}{0.45\textwidth}
      \centering
      \begin{tikzpicture}
        \fill[muted!18] (0.15,-1.4) rectangle (0.85,1.4);
        \draw[space] (0,0) ellipse (2.3 and 1.3);
        \node[lab] at (2.2,1.3) {$\X$};
        \foreach \x in {-1.6,-0.6,0.15,0.85,1.7} {
          \pgfmathsetmacro{\h}{1.3*sqrt(1-(\x/2.3)^2)}
          \draw[strip] (\x,-\h) -- (\x,\h);
        }
        \node[note] at (0.5,-1.85) {$p(x) = v$\\(constant on every strip)};
      \end{tikzpicture}
    \end{column}
    \begin{column}{0.53\textwidth}
      On every strip,
      \[
        \E\big[\,c(x)\,(y - v)\,\big] \approx 0 .
      \]

      To functions in $\C$, on a strip, $(x,y)$ looks like $(x,\tilde y)$
      where $\tilde y \sim \Ber(v)$.
    \end{column}
  \end{columns}

  \vspace{1.2em}
  \alert{A computational analogue of learning up to irreducible noise.}
\end{frame}

\begin{frame}{How to construct such predictors}
  Start with any predictor $p$.

  \textcolor{muted}{For fairness applications, start with a $p$ that has
  high accuracy / low loss.}

  \vspace{0.8em}
  Make $p$ multiaccurate / multicalibrated \alert{without increasing}
  \[
    \E\big[\,(p(x) - y)^2\,\big] .
  \]
\end{frame}

\begin{frame}{Weak agnostic learning}
  $\D$ a distribution over $\X \times \{0,1\}$; \ $\C$ a class of functions.

  \vspace{0.6em}
  \alert{$(\alpha,\beta)$-weak agnostic learner}, for $\alpha \ge \beta > 0$:
  \begin{itemize}
    \item if there is a $c \in \C$ with $\err(c) \le \tfrac12 - \alpha$,
    \item then output $h : \X \to \{0,1\}$ with $\err(h) \le \tfrac12 - \beta$.
  \end{itemize}

  \pause
  \vspace{0.6em}
  \textbf{Theorem} \textcolor{muted}{(Kalai--Kanade '09; Feldman '10)}\\
  Weak agnostic learning is equivalent to strong agnostic learning,
  even for a fixed marginal distribution over $\X$.
\end{frame}

\begin{frame}{Auditing for multiaccuracy violations}
  If $p : \X \to [0,1]$ is not multiaccurate, there is a $c \in \C$ with
  \[
    \big|\,\E\big[\,c(x)\,(y - p(x))\,\big]\,\big| \ \text{noticeable.}
  \]

  Use a weak agnostic learner to find $h$.

  \pause
  Can show that
  \[
    \E\Big[\big(\Pi_{[0,1]}\big(p(x) + \eta\, h(x)\big) - y\big)^2\Big]
    \ \text{is noticeably smaller than} \
    \E\big[\,(p(x) - y)^2\,\big] .
  \]
  \textcolor{muted}{$\Pi_{[0,1]}(f)$ clips the range of $f$ to $[0,1]$.}
\end{frame}

\begin{frame}{Multigroup fairness \& learning}
  Multicalibration with $\C$ gives a strong form of agnostic learning
  \textcolor{muted}{(also called swap agnostic learning).}

  \pause
  \vspace{1.5em}
  \alert{What about just multiaccuracy?}
\end{frame}
